\ifdefined\gkver\else\def\gkver{student}\fi
\documentclass[\gkver]{gaokaozhenti}
\begin{document}

\chapter{2009年山东理综}
%% paperType: 理综物理部分

\begin{enumerate}
\item
%% number: 16
%% paperName: 2009年山东理综第 16 题 4 分
%% typeId: 1
%% type: 单选题
%% chapter: 力物体的平衡
%% point: 三力平衡
%% method: 无
%% score: 4
%% degree: 650
%% duplicateId: 0
%% body:
如图所示，光滑半球形容器固定在水平面上，O 为球心，一质量为 $m$ 的小滑块，在水平力 $F$ 的作用下静止 P 点。设滑块所受支持力为 $F_{N}$。OF 与水平方向的夹角为 $\theta$。下列关系正确的是\xzanswer{A}

\onepicture[w=5.0cm]{figs/16.png}

\fourchoices[answer=A]
{$F=\frac{mg}{\tan\theta}$}
{$F=mg\tan\theta$}
{$F_{N}=\frac{mg}{\tan\theta}$}
{$F_{N}=mg\tan\theta$}

%% answer: A
%% memo:
\memoanswer{
本题考查受力分析。对小滑块受力分析如图所示，根据三角形定则可得 $F=\frac{mg}{\tan\theta}$，$F_{N}=\frac{mg}{\sin\theta}$，所以 A 正确。

提示：支持力的方向垂直于接触面，即指向圆心。正交分解列式求解也可。

\onepicture[w=4.5cm]{figs/16_an.png}
}

\item
%% number: 17
%% paperName: 2009年山东理综第 17 题 4 分
%% typeId: 1
%% type: 单选题
%% chapter: 运动和力
%% point: 牛顿定律的应用
%% method: 无
%% score: 4
%% degree: 650
%% duplicateId: 0
%% body:
某物体做直线运动的 $v-t$ 图象如图甲所示，据此判断图乙（$F$ 表示物体所受合力，$x$ 表示物体的位移）四个选项中正确的是\xzanswer{B}

\onepicture[num=甲, w=4.5cm]{figs/17a.png}

\fourchoices[answer=B, ispicture=true, h=2.2cm]
{figs/17b.png}
{figs/17c.png}
{figs/17d.png}
{figs/17e.png}

%% answer: B
%% memo:
\memoanswer{
由图甲可知前两秒物体做初速度为零的匀加速直线运动，所以前两秒受力恒定；2s$\sim$4s 做正方向匀加速直线运动，所以受力为负，且恒定；4s$\sim$6s 做负方向匀加速直线运动，所以受力为负，且恒定；6s$\sim$8s 做负方向匀减速直线运动，所以受力为正，且恒定。综上分析 B 正确。

提示：在 $v-t$ 图象中倾斜的直线表示物体做匀变速直线运动，加速度恒定，受力恒定。
}

\item
%% number: 18
%% paperName: 2009年山东理综第 18 题 4 分
%% typeId: 2
%% type: 多选题
%% chapter: 曲线运动
%% point: 人造地球卫星
%% method: 无
%% score: 4
%% degree: 450
%% duplicateId: 0
%% body:
2008 年 9 月 25 日至 28 日我国成功实施了“神舟”七号载人航天飞行并实现了航天员首次出舱。飞船先沿椭圆轨道飞行，后在远地点 343 千米处点火加速，由椭圆轨道变成高度为 343 千米的圆轨道，在此圆轨道上飞船运行周期约为 90 分钟。下列判断正确的是\xzanswer{BC}

\fourchoices[answer=BC]
{飞船变轨前后的机械能相等}
{飞船在圆轨道上时航天员出舱前后都处于失重状态}
{飞船在此圆轨道上运动的角速度大于同步卫星运动的角速度}
{飞船变轨前通过椭圆轨道远地点时的加速度大于变轨后沿圆轨道运动的加速度}

%% answer: BC
%% memo:
\memoanswer{
飞船点火变轨，前后的机械能不守恒，所以 A 不正确。飞船在圆轨道上时万有引力来提供向心力，航天员出舱前后都处于失重状态，B 正确。飞船在此圆轨道上运动的周期 90 分钟小于同步卫星运动的周期 24 小时，根据 $T=\frac{2\pi}{\omega}$ 可知，飞船在此圆轨道上运动的角速度大于同步卫星运动的角速度，C 正确。飞船变轨前通过椭圆轨道远地点时只有万有引力来提供加速度，变轨后沿圆轨道运动也是只有万有引力来提供加速度，所以相等，D 不正确。

提示：若物体除了重力、弹性力做功以外，还有其他力（非重力、弹性力）不做功，且其他力做功之和不为零，则机械能不守恒。根据万有引力等于卫星做圆周运动的向心力可求卫星的速度、周期、动能、动量等状态量。由 $G\frac{Mm}{r^{2}}=m\frac{v^{2}}{r}$ 得 $v=\sqrt{\frac{GM}{r}}$，由 $G\frac{Mm}{r^{2}}=m\left(\frac{2\pi}{T}\right)^{2}r$ 得 $T=2\pi\sqrt{\frac{r^{3}}{GM}}$，由 $G\frac{Mm}{r^{2}}=m\omega^{2}r$ 得 $\omega=\sqrt{\frac{GM}{r^{3}}}$，$G\frac{Mm}{r^{2}}=ma_{n}$ 可求向心加速度。
}

\item
%% number: 19
%% paperName: 2009年山东理综第 19 题 4 分
%% typeId: 2
%% type: 多选题
%% chapter: 交流电
%% point: 变压器
%% method: 无
%% score: 4
%% degree: 450
%% duplicateId: 0
%% body:
某小型水电站的电能输送示意图如下。发电机的输出电压为 $200\UV$，输电线总电阻为 $r$，升压变压器原副线圈匝数分别为 $n_{1}$、$n_{2}$。降压变压器原副线圈匝数分别为 $n_{3}$、$n_{4}$（变压器均为理想变压器）。要使额定电压为 $220\UV$ 的用电器正常工作，则\xzanswer{AD}

\onepicture[w=7.0cm]{figs/19.png}

\fourchoices[answer=AD]
{$\frac{n_{2}}{n_{1}}>\frac{n_{3}}{n_{4}}$}
{$\frac{n_{2}}{n_{1}}<\frac{n_{3}}{n_{4}}$}
{升压变压器的输出电压等于降压变压器的输入电压}
{升压变压器的输出功率大于降压变压器的输入功率}

%% answer: AD
%% memo:
\memoanswer{
根据变压器工作原理可知 $\frac{n_{1}}{n_{2}}=\frac{220}{U_{2}}$，$\frac{n_{3}}{n_{4}}=\frac{U_{3}}{220}$，由于输电线上损失一部分电压，升压变压器的输出电压大于降压变压器的输入电压，有 $U_{2}>U_{3}$，所以 $\frac{n_{2}}{n_{1}}>\frac{n_{3}}{n_{4}}$，A 正确，BC 不正确。升压变压器的输出功率等于降压变压器的输入功率加上输电线损失功率，D 正确。

提示：理想变压器的两个基本公式是：（1）$\frac{U_{1}}{U_{2}}=\frac{n_{1}}{n_{2}}$，即对同一变压器的任意两个线圈，都有电压和匝数成正比。（2）$P_{1}=P_{2}$，即无论有几个副线圈在工作，变压器的输入功率总等于所有输出功率之和。只有当变压器只有一个副线圈工作时，才有 $U_{1}I_{1}=U_{2}I_{2}$，$\frac{I_{1}}{I_{2}}=\frac{n_{2}}{n_{1}}$。

远距离输电，从图中应该看出功率之间的关系是：$P_{1}=P_{2}$，$P_{3}=P_{4}$，$P_{1}'=P_{r}=P_{2}$。电压之间的关系是：$\frac{U_{1}}{U_{2}}=\frac{n_{1}}{n_{2}}$，$\frac{U_{3}}{U_{4}}=\frac{n_{3}}{n_{4}}$，$U_{2}=U_{r}+U_{3}$。电流之间的关系是：$\frac{I_{1}}{I_{2}}=\frac{n_{2}}{n_{1}}$，$\frac{I_{3}}{I_{4}}=\frac{n_{4}}{n_{3}}$，$I_{2}=I_{r}=I_{3}$。输电线上的功率损失和电压损失也是需要特别注意的。分析和计算时都必须用 $P_{r}=I_{2}^{2}r$，$U_{r}=I_{2}r$，而不能用 $P_{r}=\frac{U_{r}^{2}}{r}$。特别重要的是要会分析输电线上的功率损失 $P_{r}=\left(\frac{P_{1}}{U_{2}}\right)^{2}\cdot\rho\frac{l}{S}\propto\frac{1}{U_{2}^{2}S}$。
}

\item
%% number: 20
%% paperName: 2009年山东理综第 20 题 4 分
%% typeId: 2
%% type: 多选题
%% chapter: 电场
%% point: 电场线
%% method: 无
%% score: 4
%% degree: 450
%% duplicateId: 0
%% body:
如图所示，在 $x$ 轴上关于原点 O 对称的两点固定放置等量异种点电荷 $+Q$ 和 $-Q$，$x$ 轴上的 P 点位于 $-Q$ 的右侧。下列判断正确的是\xzanswer{AC}

\onepicture[w=6.0cm]{figs/20.png}

\fourchoices[answer=AC]
{在 $x$ 轴上还有一点与 P 点电场强度相同}
{在 $x$ 轴上还有两点与 P 点电场强度相同}
{若将一试探电荷 $+q$ 从 P 点移至 O 点，电势能增大}
{若将一试探电荷 $+q$ 从 P 点移至 O 点，电势能减小}

%% answer: AC
%% memo:
\memoanswer{
根据等量正负点电荷的电场分布可知，在 $x$ 轴上还有一点与 P 点电场强度相同，即和 P 点关于 O 点对称，A 正确。若将一试探电荷 $+q$ 从 P 点移至 O 点，电场力先做正功后做负功，所以电势能先减小后增大。一般规定无穷远电势为零，过 O 点的中垂线电势也为零，所以试探电荷 $+q$ 在 P 点时电势能为负值，移至 O 点时电势能为零，所以电势能增大，C 正确。

提示：熟悉掌握等量正负点电荷的电场分布。知道 $W_{AB}=E_{PA}-E_{PB}$，即电场力做正功，电势能转化为其他形式的能，电势能减少；电场力做负功，其他形式的能转化为电势能，电势能增加，即 $W=-\Delta E$。

\onepicture[w=5.0cm]{figs/20_an.png}
}

\item
%% number: 21
%% paperName: 2009年山东理综第 21 题 4 分
%% typeId: 3
%% type: 填空题
%% chapter: 电磁感应
%% point: 切割磁感线电动势
%% method: 无
%% score: 4
%% degree: 450
%% duplicateId: 0
%% body:
如图所示，一导线弯成半径为 $a$ 的半圆形闭合回路。虚线 MN 右侧有磁感应强度为 $B$ 的匀强磁场，方向垂直于回路所在的平面。回路以速度 $v$ 向右匀速进入磁场，直径 CD 始终与 MN 垂直。从 D 点到达边界开始到 C 点进入磁场为止，下列结论正确的是\xzanswer{ACD}

\onepicture[w=5.0cm, align=right]{figs/21.png}

\fourchoices[answer=ACD]
{感应电流方向不变}
{CD 段直线始终不受安培力}
{感应电动势最大值 $E=Bav$}
{感应电动势平均值 $\overline{E}=\frac{1}{4}\pi Bav$}

%% answer: ACD
\jdanswer{ACD}
%% memo:
\memoanswer{
在闭合电路进入磁场的过程中，通过闭合电路的磁通量逐渐增大，根据楞次定律可知感应电流的方向为逆时针方向不变，A 正确。根据左手定则可以判断，受安培力向下，B 不正确。当半圆闭合回路进入磁场一半时，即这时等效长度最大为 $a$，这时感应电动势最大 $E=Bav$，C 正确。感应电动势平均值 $\overline{E}=\frac{\Delta\phi}{\Delta t}=\frac{B\cdot\frac{1}{2}\pi a^{2}}{\frac{2a}{v}}=\frac{1}{4}\pi Bav$，D 正确。

提示：感应电动势公式 $\overline{E}=\frac{\Delta\phi}{\Delta t}$ 只能来计算平均值，利用感应电动势公式 $E=Blv$ 计算时，$l$ 应是等效长度，即垂直切割磁感线的长度。
}

\item
%% number: 22
%% paperName: 2009年山东理综第 22 题 4 分
%% typeId: 2
%% type: 多选题
%% chapter: 功和能
%% point: 功能原理
%% method: 无
%% score: 4
%% degree: 450
%% duplicateId: 0
%% body:
图示为某探究活动小组设计的节能运动系统。斜面轨道倾角为 $30^\circ$，质量为 $M$ 的木箱与轨道的动摩擦因数为 $\frac{\sqrt{3}}{6}$。木箱在轨道端时，自动装货装置将质量为 $m$ 的货物装入木箱，然后木箱载着货物沿轨道无初速滑下，与轻弹簧被压缩至最短时，自动卸货装置立刻将货物卸下，然后木箱恰好被弹回到轨道顶端，再重复上述过程。下列选项正确的是\xzanswer{BC}

\onepicture[w=4.5cm]{figs/22.png}

\fourchoices[answer=BC]
{$m=M$}
{$m=2M$}
{木箱不与弹簧接触时，上滑的加速度大于下滑的加速度}
{在木箱与货物从顶端滑到最低点的过程中，减少的重力势能全部转化为弹簧的弹性势能}

%% answer: BC
%% memo:
\memoanswer{
受力分析可知，下滑时加速度为 $g\sin\theta-\mu g\cos\theta$，上滑时加速度为 $g\sin\theta+\mu g\cos\theta$，所以 C 正确。设下滑的距离为 $l$，根据能量守恒有 $\mu(m+M)gl\cos\theta+\mu Mgl\cos\theta=mgl\sin\theta$，得 $m=2M$。也可以根据除了重力、弹性力做功以外，还有其他力（非重力、弹性力）做的功之和等于系统机械能的变化量，B 正确。在木箱与货物从顶端滑到最低点的过程中，减少的重力势能转化为弹簧的弹性势能和内能，所以 D 不正确。

提示：能量守恒定律的理解及应用。
}

\item
%% number: 23
%% paperName: 2009年山东理综第 23 题 7 分
%% typeId: 4
%% type: 实验题
%% chapter: 电路
%% point: 传感器
%% method: 无
%% score: 7
%% degree: 450
%% duplicateId: 0
%% body:
请完成以下两小题。

\textbf{（1）}某同学在家中尝试验证平行四边形定则，他找到三条相同的橡皮筋（遵循胡克定律）和若干小事物，以及刻度尺、三角板、铅笔、细绳、白纸、钉子，设计了如下实验：将两条橡皮筋的一端分别系在墙上的两个钉子 A、B 上，另一端与第三条橡皮筋连接，结点为 O，将第三条橡皮筋的另一端通过细绳挂一重物。

\onepicture[w=3.2cm, align=right]{figs/23.png}

①为完成实验，下述操作中必需的是\tkanswer[3cm]{}。

（A）测量细绳的长度　　（B）测量橡皮筋的原长

（C）测量悬挂重物后橡皮筋的长度　　（D）记录悬挂重物后结点 O 的位置

②钉子位置固定，欲利用现有器材，改变条件再次验证，可采用的方法是\tkanswer[3cm]{}。

\textbf{（2）}为了节能和环保，一些公共场所使用光控开关控制照明系统。光控开关可采用光敏电阻来控制，光敏电阻是阻值随着光的照度而发生变化的元件（照度可以反映光的强弱，光越强照度越大，照度单位为 $\Ulx$）。某光敏电阻 $R_{P}$ 在不同照度下的阻值如下表：

\begin{center}
\begin{tabular}{|c|c|c|c|c|c|c|}
\hline
照度（$\Ulx$） & 0.2 & 0.4 & 0.6 & 0.8 & 1.0 & 1.2 \\
\hline
电阻（$\UkO$） & 75 & 40 & 28 & 23 & 20 & 18 \\
\hline
\end{tabular}
\end{center}

①根据表中数据，请在给定的坐标系中描绘出阻值随照度变化的曲线，并说明阻值随照度变化的特点。

\onepicture[w=6.5cm]{figs/23a.png}

②如图所示，当 1、2 两端所加电压上升至 $2\UV$ 时，控制开关自动启动照明系统，请利用下列器材设计一个简单电路，给 1、2 两端提供电压，要求当天色渐暗照度降低至 $1.0\Ulx$ 时启动照明系统，在虚线框内完成电路原理图。（不考虑控制开关对所设计电路的影响）

\onepicture[w=5.0cm]{figs/23b.png}

提供的器材如下：

光敏电阻 $R_{P}$（符号如图所示，阻值见上表）

\onepicture[w=1.3cm, align=right]{figs/23c.png}

直流电源 $E$（电动势 $3\UV$，内阻不计）；

定值电阻：$R_{1}=10\UkO$，$R_{2}=20\UkO$，$R_{3}=40\UkO$（限选其中之一并在图中标出）；开关 S 及导线若干。

%% answer:
\jdanswer{（1）① BCD；② 更换不同的小重物。（2）① 光敏电阻的阻值随光照变化的曲线如解析图所示；特点：光敏电阻的阻值随光照强度的增大非线性减小。② 电路原理图如解析图所示，此时光敏电阻阻值为 $20\UkO$，分压电阻选用 $R_{1}$。}
%% memo:
\memoanswer{
（1）本题原卷及材料中未提供详解，待补充。

（2）当控制开关自动启动照明系统时，给 1、2 两端提供电压，要求当天色渐暗照度降低至 $1.0\Ulx$ 时启动照明系统，即此时光敏电阻阻值为 $20\UkO$，两端电压为 $2\UV$，电源电动势为 $3\UV$，所以应加上一个分压电阻，分压电阻阻值为 $10\UkO$，即选用 $R_{1}$。

\onepicture[w=6.5cm]{figs/23_an1.png}

\onepicture[w=5.0cm]{figs/23_an2.png}
}

\item
%% number: 24
%% paperName: 2009年山东理综第 24 题 15 分
%% typeId: 6
%% type: 计算题
%% chapter: 运动和力
%% point: 传送带
%% method: 无
%% score: 15
%% degree: 450
%% duplicateId: 0
%% body:
如图所示，某货场将质量为 $m_{1}=100\Ukg$ 的货物（可视为质点）从高处运送至地面，为避免货物与地面发生撞击，现利用固定于地面的光滑四分之一圆轨道，使货物从轨道顶端无初速滑下，轨道半径 $R=1.8\Um$。地面上紧靠轨道依次排放两块完全相同的木板 A、B，长度均为 $l=2\Um$，质量均为 $m_{2}=100\Ukg$，木板上表面与轨道末端相切。货物与木板间的动摩擦因数为 $\mu_{1}$，木板与地面间的动摩擦因数 $\mu_{2}=0.2$。（最大静摩擦力与滑动摩擦力大小相等，取 $g=10\Umsq$）

\begin{enumerate}
	\item
	求货物到达圆轨道末端时对轨道的压力。

	\item
	若货物滑上木板 A 时，木板不动，而滑上木板 B 时，木板 B 开始滑动，求 $\mu_{1}$ 应满足的条件。

	\item
	若 $\mu_{1}=0.5$，求货物滑到木板 A 末端时的速度和在木板 A 上运动的时间。
\end{enumerate}
\onepicture[w=7.0cm, align=right]{figs/24.png}

%% answer:
\jdanswer{
\begin{enumerate}
\item $3000\UN$
\item $0.4<\mu_{1}\leq0.6$
\item $4\Ums$；$0.4\Us$
\end{enumerate}
}
%% memo:
\memoanswer{
（1）设货物滑到圆轨道末端时的速度为 $v_{0}$，对货物下滑过程根据机械能守恒定律得 $m_{1}gR=\frac{1}{2}m_{1}v_{0}^{2}$ ①，设货物在轨道末端所受支持力的大小为 $F_{N}$，根据牛顿第二定律得 $F_{N}-m_{1}g=m_{1}\frac{v_{0}^{2}}{R}$ ②，联立以上两式代入数据得 $F_{N}=3000\UN$ ③。根据牛顿第三定律，货物到达圆轨道末端时对轨道的压力大小为 $3000\UN$，方向竖直向下。

（2）若滑上木板 A 时，木板不动，由受力分析得 $\mu_{1}m_{1}g\leq\mu_{2}(m_{1}+2m_{2})g$ ④；若滑上木板 B 时，木板 B 开始滑动，由受力分析得 $\mu_{1}m_{1}g>\mu_{2}(m_{1}+m_{2})g$ ⑤。联立④⑤式代入数据得 $0.4<\mu_{1}\leq0.6$ ⑥。

（3）$\mu_{1}=0.5$，由⑥式可知，货物在木板 A 上滑动时，木板不动。设货物在木板 A 上做减速运动时的加速度大小为 $a_{1}$，由牛顿第二定律得 $\mu_{1}m_{1}g=m_{1}a_{1}$ ⑦，设货物滑到木板 A 末端时的速度为 $v_{1}$，由运动学公式得 $v_{1}^{2}-v_{0}^{2}=-2a_{1}l$ ⑧，联立①⑦⑧式代入数据得 $v_{1}=4\Ums$ ⑨，设在木板 A 上运动的时间为 $t$，由运动学公式得 $v_{1}=v_{0}-a_{1}t$ ⑩，联立①⑦⑨⑩式代入数据得 $t=0.4\Us$。
}

\item
%% number: 25
%% paperName: 2009年山东理综第 25 题 18 分
%% typeId: 6
%% type: 计算题
%% chapter: 磁场
%% point: 带电粒子在复合场中的运动
%% method: 无
%% score: 18
%% degree: 250
%% duplicateId: 0
%% body:
如图甲所示，建立 $Oxy$ 坐标系，两平行极板 P、Q 垂直于 $y$ 轴且关于 $x$ 轴对称，极板长度和板间距均为 $l$，第一、四象限有磁场，方向垂直于 $Oxy$ 平面向里。位于极板左侧的粒子源沿 $x$ 轴向右连续发射质量为 $m$、电荷量为 $+q$、速度相同、重力不计的带电粒子，在 $0\sim3t_{0}$ 时间内两板间加上如图乙所示的电压（不考虑极板边缘的影响）。已知 $t=0$ 时刻进入两板间的带电粒子恰好在 $t_{0}$ 时刻经极板边缘射入磁场。上述 $m$、$q$、$l$、$t_{0}$、$B$ 为已知量。（不考虑粒子间相互影响及返回板间的情况）

\begin{enumerate}
	\item
	求电压 $U$ 的大小。

	\item
	求 $\frac{t_{0}}{2}$ 时进入两板间的带电粒子在磁场中做圆周运动的半径。

	\item
	何时把两板间的带电粒子在磁场中的运动时间最短？求此最短时间。
\end{enumerate}
\twopicture[numA=甲, numB=乙, labelprefix={图}, w=6.5cm, align=right]
{figs/25a.png}
{figs/25b.png}

%% answer:
\jdanswer{
\begin{enumerate}
\item $U_{0}=\frac{ml^{2}}{qt_{0}^{2}}$
\item $R=\frac{\sqrt{5}ml}{2qBt_{0}}$
\item $2t_{0}$ 时刻进入的带电粒子在磁场中运动时间最短，$t_{\min}=\frac{\pi m}{2qB}$
\end{enumerate}
}
%% memo:
\memoanswer{
（1）$t=0$ 时刻进入两极板的带电粒子在电场中做匀变速曲线运动，$t_{0}$ 时刻刚好从极板边缘射出，在 $y$ 轴负方向偏移的距离为 $\frac{l}{2}$，则有 $E=\frac{U_{0}}{l}$ ①，$Eq=ma$ ②，$\frac{l}{2}=\frac{1}{2}at_{0}^{2}$ ③，联立以上三式，解得两极板间偏转电压为 $U_{0}=\frac{ml^{2}}{qt_{0}^{2}}$ ④。

（2）$\frac{t_{0}}{2}$ 时刻进入两极板的带电粒子，前 $\frac{t_{0}}{2}$ 时间在电场中偏转，后 $\frac{t_{0}}{2}$ 时间两极板没有电场，带电粒子做匀速直线运动。带电粒子沿 $x$ 轴方向的分速度大小为 $v_{0}=\frac{l}{t_{0}}$ ⑤，带电粒子离开电场时沿 $y$ 轴负方向的分速度大小为 $v_{y}=a\cdot\frac{1}{2}t_{0}$ ⑥，带电粒子离开电场时的速度大小为 $v=\sqrt{v_{x}^{2}+v_{y}^{2}}$ ⑦。设带电粒子离开电场进入磁场做匀速圆周运动的半径为 $R$，则有 $Bqv=m\frac{v^{2}}{R}$ ⑧，联立③⑤⑥⑦⑧式解得 $R=\frac{\sqrt{5}ml}{2qBt_{0}}$ ⑨。

（3）$2t_{0}$ 时刻进入两极板的带电粒子在磁场中运动时间最短。带电粒子离开电场时沿 $y$ 轴正方向的分速度为 $v_{y}'=at_{0}$ ⑩，设带电粒子离开电场时速度方向与 $y$ 轴正方向的夹角为 $\alpha$，则 $\tan\alpha=\frac{v_{0}}{v_{y}'}$，联立③⑤⑩式解得 $\alpha=\frac{\pi}{4}$，带电粒子在磁场运动的轨迹图如图所示，圆弧所对的圆心角为 $2\alpha=\frac{\pi}{2}$，所求最短时间为 $t_{\min}=\frac{1}{4}T$，带电粒子在磁场中运动的周期为 $T=\frac{2\pi m}{qB}$，联立以上两式解得 $t_{\min}=\frac{\pi m}{2qB}$。

\onepicture[w=6.5cm]{figs/25_an.png}
}

\item
%% number: 36
%% paperName: 2009年山东理综第 36 题 8 分
%% typeId: 6
%% type: 计算题
%% chapter: 气体性质
%% point: 盖吕萨克定律
%% method: 无
%% score: 8
%% degree: 650
%% duplicateId: 0
%% body:
一定质量的理想气体由状态 A 经状态 B 变为状态 C，其中 A$\to$B 过程为等压变化，B$\to$C 过程为等容变化。已知 $V_{A}=0.3\Umc$，$T_{A}=T_{C}=300\UK$、$T_{B}=400\UK$。

\begin{enumerate}
	\item
	求气体在状态 B 时的体积。

	\item
	说明 B$\to$C 过程压强变化的微观原因。

	\item
	设 A$\to$B 过程气体吸收热量为 $Q_{1}$，B$\to$C 过程气体放出热量为 $Q_{2}$，比较 $Q_{1}$、$Q_{2}$ 的大小并说明原因。
\end{enumerate}

%% answer:
\jdanswer{
\begin{enumerate}
\item $V_{B}=0.4\Umc$
\item 气体体积不变，分子密集程度不变，温度变小，气体分子平均动能减小，导致气体压强减小。
\item $Q_{1}>Q_{2}$
\end{enumerate}
}
%% memo:
\memoanswer{
（1）设气体在 B 状态时的体积为 $V_{B}$，由盖—吕萨克定律得 $\frac{V_{A}}{T_{A}}=\frac{V_{B}}{T_{B}}$，代入数据得 $V_{B}=0.4\Umc$。

（2）微观原因：气体体积不变，分子密集程度不变，温度变小，气体分子平均动能减小，导致气体压强减小。

（3）$Q_{1}$ 大于 $Q_{2}$。因为 $T_{A}=T_{C}$，故 A$\to$B 增加的内能与 B$\to$C 减小的内能相同，而 A$\to$B 过程气体对外做正功，B$\to$C 过程气体不做功，由热力学第一定律可知 $Q_{1}$ 大于 $Q_{2}$。
}

\item
%% number: 37
%% paperName: 2009年山东理综第 37 题 4 分
%% typeId: 6
%% type: 计算题
%% chapter: 光学
%% point: 全反射
%% method: 无
%% score: 4
%% degree: 650
%% duplicateId: 0
%% body:
请完成以下两小题。

\textbf{（1）}如图所示为一列简谐波在 $t=0$ 时的波形图，介质中的质点 P 做简谐运动的表达式为 $y=4\sin5\pi t$，求该波的速度，并画出 $t=0.3\Us$ 时的波形图（至少画出一个波长）。

\onepicture[w=6.0cm]{figs/37a.png}

\textbf{（2）}一束单色光由左侧射入装有清水的薄壁圆柱形玻璃杯，如图所示为过轴线的截面图，调整入射角 $\alpha$，光线恰好在水和空气的界面上发生全反射，已知水对此光的折射率为 $\frac{4}{3}$，求 $\alpha$ 的值。

\onepicture[w=4.0cm, align=right]{figs/37b.png}

%% answer:
\jdanswer{
\begin{enumerate}
\item $v=10\Ums$，$t=0.3\Us$ 时的波形图如解析图所示
\item $\sin\alpha=\frac{\sqrt{7}}{3}$
\end{enumerate}
}
%% memo:
\memoanswer{
（1）由简谐运动的表达式可知 $\omega=5\pi\Urads$，$t=0$ 时刻质点 P 向上运动，故波沿 $x$ 轴正方向传播。由波形图读出波长 $\lambda=4\Um$，$T=\frac{2\pi}{\omega}$，由波速公式 $v=\frac{\lambda}{T}$，联立以上两式代入数据可得 $v=10\Ums$。$t=0.3\Us$ 时的波形图如图所示。

\onepicture[w=6.0cm]{figs/37_an.png}

（2）当光线在水面发生全反射时有 $\sin C=\frac{1}{n}$，当光线从左侧射入时，由折射定律有 $\frac{\sin\alpha}{\sin\left(\frac{\pi}{2}-C\right)}=n$，联立这两式代入数据可得 $\sin\alpha=\frac{\sqrt{7}}{3}$。
}

\item
%% number: 38
%% paperName: 2009年山东理综第 38 题 4 分
%% typeId: 6
%% type: 计算题
%% chapter: 运动和力
%% point: 动量守恒定律
%% method: 无
%% score: 4
%% degree: 650
%% duplicateId: 0
%% body:
请完成以下两小题。

\textbf{（1）}历史上在利用加速器实现的核反应中，是用加速后动能为 $0.5\UMeV$ 的质子 $\ce{^1_1H}$ 轰击静止的 X，生成两个动能均为 $8.9\UMeV$ 的 $\ce{^4_2He}$。（$1\UMeV=1.6\times10^{-13}\UJ$）

①上述核反应方程为\tkanswer[4cm]{}。

②质量亏损为\tkanswer[3cm]{} $\Ukg$。

\textbf{（2）}如图所示，光滑水平面轨道上有三个木块 A、B、C，质量分别为 $m_{A}=m_{C}=2m$，$m_{B}=m$，A、B 用细绳连接，中间有一压缩的弹簧（弹簧与滑块不栓接）。开始时 A、B 以共同速度 $v_{0}$ 运动，C 静止。某时刻细绳突然断开，A、B 被弹开，然后 B 又与 C 发生碰撞并粘在一起，最终三滑块速度恰好相同。求 B 与 C 碰撞前 B 的速度。

\onepicture[w=6.0cm, align=right]{figs/38.png}

%% answer:
\jdanswer{
\begin{enumerate}
\item ①$\ce{^1_1H + ^7_3Li -> ^4_2He + ^4_2He}$（或 $\ce{^1_1H + ^7_3X -> ^4_2He + ^4_2He}$）；②$3.1\times10^{-29}\Ukg$
\item $v_{B}=\frac{9}{5}v_{0}$
\end{enumerate}
}
%% memo:
\memoanswer{
（1）$\ce{^1_1H + ^7_3Li -> ^4_2He + ^4_2He}$（或 $\ce{^1_1H + ^7_3X -> ^4_2He + ^4_2He}$），$\Delta m=\frac{E}{c^{2}}=3.1\times10^{-29}\Ukg$。

（2）设共同速度为 $v$，分离后 B 的速度为 $v_{B}$，由动量守恒定律有 $(m_{A}+m_{B})v_{0}=m_{A}v+m_{B}v_{B}$，$m_{B}v_{B}=(m_{B}+m_{C})v$，联立这两式得 B 和 C 碰撞前 B 的速度为 $v_{B}=\frac{9}{5}v_{0}$。
}

\end{enumerate}

\end{document}
